% GENERATED, edit the French unit or the translation table
% GENERATED from content/data/enrichment.json by mxbuild.enrichment; do not edit.
\mxunit{enrichment-bisection-exercises}{Bracketing a root by bisection}
\mxsubtitle{Upper high school · analysis and algorithms \textperiodcentered{} Progressive exercises \textperiodcentered{} 12}

\begin{note}
\textbf{Given information.} For exercise 1: seek a zero of $f(x)=x^{2}-2$ on $\intff{1}{2}$. After two halvings, the retained interval is $\intff{\num{1,25}}{\num{1,5}}$.
\end{note}

\begin{exercise}
Carry out a third step of the example.
\end{exercise}

\begin{exercise}
Starting with width 1, how many halvings guarantee width at most $\frac{1}{1000}$?
\end{exercise}

\begin{exercise}
Why does the sign change of $f(x)=\frac{1}{x}$ between $-1$ and 1 not prove the existence of a zero?
\end{exercise}

\begin{exercise}
A root lies in $\intff{\num{1,4}}{\num{1,5}}$. What maximum error is guaranteed by midpoint $\num{1,45}$?
\end{exercise}

